\documentclass[conference]{IEEEtran}
\IEEEoverridecommandlockouts

\usepackage{cite}
\usepackage{amsmath,amssymb,amsfonts}
\usepackage{graphicx}
\usepackage{textcomp}
\usepackage{xcolor}
\usepackage{booktabs}
\usepackage[spanish,es-tabla]{babel}
\usepackage[T1]{fontenc}
\usepackage[utf8]{inputenc}
\usepackage{hyperref}
\usepackage{listings}

% Placeholder visible para las cifras que faltan completar desde
% evidence/metrics/ al cerrar las corridas oficiales. Buscar "PEND" antes de
% entregar: ninguna cifra de este documento se escribe a mano (regla 8 de
% CLAUDE.md), asi que cada PEND se reemplaza copiando el valor exacto que
% imprime compute_metrics.py.
\newcommand{\PEND}[1]{\textcolor{red}{\textbf{[PEND: #1]}}}

\def\BibTeX{{\rm B\kern-.05em{\sc i\kern-.025em b}\kern-.08em
    T\kern-.1667em\lower.7ex\hbox{E}\kern-.125emX}}

\begin{document}

\title{Persistencia indebida de informaci\'on e instrucciones en la memoria\\
conversacional de asistentes LLM: aislamiento, expiraci\'on\\
y eliminaci\'on verificable}

\author{
\IEEEauthorblockN{Isaac David Burgos Cervantes, Jose Luis Lancheros Aroya, Gina Sof\'ia Garc\'ia Zapata}
\IEEEauthorblockA{
Equipo FDSI-GP-06 \\
Seminario Aplicado de Ciberseguridad en IA, FDSI/SPTI, Escuela Colombiana de Ingenier\'ia \\
\texttt{github.com/sofiapeace/fdsi-gp06-secure-conversational-memory}
}
}

\maketitle

\begin{abstract}
Los asistentes conversacionales con memoria de largo plazo guardan hechos de
cada turno y los reinyectan en sesiones futuras, lo que convierte ese
almac\'en en un espacio de escritura permanente que de facto controla qui\'en
influye en el modelo. Construimos un laboratorio local y cerrado con un
orquestador propio sobre Llama 3.1 8B Instruct (Ollama) y dos perfiles sobre
\textbf{una sola base de c\'odigo}: \emph{Unsecure}, que reproduce la
arquitectura t\'ipica de "extraer y guardar" sin ning\'un control, y
\emph{Secure}, que agrega seis controles (C1--C6): aislamiento por espacio de
nombres en los dos canales de recuperaci\'on, separaci\'on entre datos e
instrucciones en el ensamblado del prompt, un filtro de escritura tipado,
procedencia firmada con HMAC, tiempo de vida con borrado verificable en tres
capas, y autorizaci\'on de herramientas por origen. En la l\'inea base
Unsecure (Hito 2, commit \texttt{3ada423}) una instrucci\'on inyectada en una
sesi\'on persisti\'o y se obedeci\'o en el 100\,\% (20/20) de sesiones nuevas
sin relaci\'on sem\'antica con ella, y el \'indice vectorial entreg\'o datos
de otro usuario en el 67\,\% de 100 sondeos (fuga de recuperaci\'on, M2a),
aunque el modelo solo los repiti\'o en el 49\,\% (fuga de divulgaci\'on,
M2b): la brecha de 18 puntos muestra que medir \'unicamente la respuesta del
modelo subestima la falla del control en un 27\,\% relativo. En la misma
m\'aquina y con el mismo modelo, repetir el perfil Unsecure con el c\'odigo del
Hito~3 dio exactamente las mismas cifras. Con C1, C2, C3 y C5 activos, el
perfil Secure obtuvo TPI~=~0\,\% (0/40, IC~95\,\% [0--8,76]), porque el filtro
de escritura rechaz\'o la instrucci\'on antes de guardarla; M2a~=~M2b~=~0\,\%
(0/100, IC [0--3,70]), con el dato ajeno presente en el almac\'en; 0 de 10
recuperaciones de un dato vencido (frente a 10 de 10); y borrado efectivo en
las tres capas (frente a una de tres). El percentil 95 de la ruta de memoria
(\(L_{mem}\), que excluye la generaci\'on de tokens) pas\'o de 83 a 86~ms,
frente a un presupuesto de 300~ms. Se discuten las amenazas a la validez y se
declara expl\'icitamente qu\'e parte del alcance qued\'o fuera por tiempo.
\end{abstract}

\begin{IEEEkeywords}
memoria conversacional, inyecci\'on de prompts, envenenamiento de memoria,
aislamiento multiusuario, borrado verificable, LLM, MITRE ATLAS, OWASP
\end{IEEEkeywords}

\section{Introducci\'on}

Los asistentes de chat con memoria de largo plazo extraen hechos de cada
turno, los guardan en una base de datos y en un \'indice vectorial, y los
vuelven a introducir en el \emph{prompt} de sesiones futuras. Ese mecanismo,
pensado para la utilidad, convierte la memoria en un espacio de escritura
permanente que en la pr\'actica decide qui\'en tiene influencia sobre lo que
el modelo dice despu\'es. Si no existen controles de aislamiento, procedencia,
tiempo de vida y borrado verificable, aparecen tres amenazas concretas: (i)
una instrucci\'on escrita por un atacante en una sesi\'on queda guardada como
si fuera un hecho del usuario y se reactiva sola en sesiones futuras, sin que
el atacante necesite seguir conectado; (ii) si la b\'usqueda por similitud no
filtra por due\~no del dato, informaci\'on de un usuario puede aparecer en la
sesi\'on de otro; y (iii) cuando el usuario pide que se olvide algo, el
borrado t\'ipico solo marca la fila en la base relacional y deja intactos el
embedding en el \'indice vectorial y las copias en cach\'e.

Este trabajo mide cu\'anto reducen esas tres amenazas seis controles
(C1--C6) frente a la misma aplicaci\'on sin ellos, usando \textbf{una sola
base de c\'odigo} con dos perfiles de configuraci\'on (\emph{unsecure} y
\emph{secure}), de forma que el \'unico diff entre arquitecturas sea el
conjunto de controles y no una reimplementaci\'on paralela que invalidar\'ia
la comparaci\'on. El laboratorio es local, cerrado, usa datos inventados
(tarjeta de prueba est\'andar de Visa \texttt{4111 1111 1111 1111}, direcciones
ficticias) y no ataca ning\'un sistema de terceros.

\textbf{Contribuciones.} (1) Un laboratorio reproducible de extremo a extremo,
con pruebas T01--T06
que aislan cada control midiendo \'unicamente lo que ese control dice medir.
(2) Una separaci\'on metodol\'ogica entre fuga de \emph{recuperaci\'on} (M2a,
lo que el \'indice entreg\'o) y fuga de \emph{divulgaci\'on} (M2b, lo que el
modelo repiti\'o), que evita atribuirle a un control el m\'erito de que el
modelo haya decidido callarse. (3) Una prueba de borrado que usa HMAC sobre el
contenido en vez de un hash plano, para que el log de auditor\'ia que
demuestra el borrado no termine siendo \'el mismo una copia recuperable del
dato borrado. (4) Evidencia reproducible versionada (\texttt{git\_sha}, semilla,
digest del modelo) para cada corrida.

\section{Trabajo relacionado}

El envenenamiento de memoria en agentes LLM es una categor\'ia de riesgo
reconocida por los tres marcos normativos que seguimos. El OWASP Top 10 for
LLM Applications 2025 cataloga la inyecci\'on de \emph{prompts} (LLM01) y la
divulgaci\'on de informaci\'on sensible (LLM02) \cite{owasp-llm10-2025}. El
OWASP Top 10 for Agentic Applications a\~nade un riesgo propio para memoria y
contexto, ASI06 \emph{Memory \& Context Poisoning} \cite{owasp-agentic-2026}.
MITRE ATLAS define una t\'ecnica dedicada, \textbf{AML.T0080 ``AI Agent Context
Poisoning''}, con la subt\'ecnica \textbf{AML.T0080.000 ``Memory''} dentro de
la t\'actica \emph{Persistence}: manipular la memoria de un LLM para persistir
cambios hacia sesiones futuras, con una mitigaci\'on dedicada,
\textbf{AML.M0031 Memory Hardening} \cite{mitre-atlas}. La fuga entre
conversaciones corresponde a \textbf{AML.T0057 LLM Data Leakage} (t\'actica
\emph{Exfiltration}) y a la recolecci\'on de datos de bases RAG,
\textbf{AML.T0085.000} (t\'actica \emph{Collection}).

Dong et al.\ demuestran MINJA, un ataque de inyecci\'on de memoria que opera
\'unicamente mediante consultas normales al agente, sin acceso privilegiado al
sistema \cite{minja-2025}. Devarangadi Sunil et al.\ proponen y calibran dos
defensas --moderaci\'on basada en confianza y saneamiento de memoria con
filtrado temporal-- que se corresponden directamente con nuestros controles
C3 y C5, y reportan un hallazgo que respalda de forma independiente una
decisi\'on de dise\~no de este laboratorio: bajo condiciones realistas, con
memoria leg\'itima de fondo ya presente, la efectividad del ataque cae
dr\'asticamente frente a un almac\'en vac\'io \cite{memory-poisoning-defense-2026}.
Evaluaron, entre otros modelos, Llama-3.1-8B-Instruct, el mismo modelo que se
usa aqu\'i. NIST clasifica estos ataques dentro de su taxonom\'ia de
\emph{adversarial machine learning} \cite{nist-aml-2025}. Ninguno de estos
trabajos mide, sobre la misma base de c\'odigo, el costo en latencia y
utilidad de desplegar los controles junto a una l\'inea base insegura
id\'entica en todo lo dem\'as: es la brecha que este laboratorio cierra.

\section{Modelo de amenaza y arquitectura}

\textbf{Activo.} El almac\'en de memoria (perfil del usuario, resumenes,
embeddings, historial) y el contexto que el orquestador ensambla para cada
sesi\'on nueva.

\textbf{Adversario.} Un usuario leg\'itimo del propio sistema (sin acceso
privilegiado) que escribe texto en un turno normal de conversaci\'on, o que
sube un documento que el sistema indexa. No se asume acceso a la
infraestructura ni a los pesos del modelo.

\subsection{Arquitectura Unsecure (Fig.~1)}

El orquestador tiene un \emph{Memory Writer} que detecta si un turno ``vale
la pena recordarlo'' con pistas de \textbf{utilidad}, no de seguridad
(``recuerda'', ``siempre'', ``prefiero'', ``a partir de ahora''), y lo
persiste sin tipar ni validar. El \emph{Memory Retriever} tiene \textbf{dos
canales}: un canal de \emph{perfil fijado}, que entrega en todas las sesiones
lo que el writer marc\'o como preferencia permanente sin pasar por ninguna
b\'usqueda, y un canal de \emph{b\'usqueda por similitud} (top-$k$ sobre
ChromaDB). Ninguno de los dos filtra por \texttt{tenant\_id} ni
\texttt{user\_id}: el dato \emph{s\'i} queda etiquetado con su due\~no en los
metadatos del \'indice vectorial, pero ninguna consulta usa esa etiqueta
como filtro. El ensamblador de \emph{prompt} concatena lo recuperado dentro
del mensaje de \textbf{sistema}, con la misma autoridad que las reglas del
operador. El almac\'en (SQLite + ChromaDB + cach\'e de sesi\'on) no tiene
tiempo de vida, ni firma de procedencia, ni borrado real: borrar una fila es
un borrado l\'ogico que deja vivo el embedding.

Cada prueba explota un canal distinto: T01 usa el perfil fijado para
demostrar persistencia sin relaci\'on sem\'antica con la pregunta; T02
usa la b\'usqueda por similitud para demostrar fuga entre usuarios. Las dos
pruebas siembran memoria leg\'itima de fondo antes de atacar (20 registros sembrados, de los que el writer guarda 15 en los dos perfiles),
porque sin ese ruido el almac\'en quedar\'ia casi vac\'io y el top-$k$
devolver\'ia el dato del ataque por simple descarte, no por la debilidad real
que se quiere demostrar.

\subsection{Arquitectura Secure (Fig.~2)}

Misma base de c\'odigo, perfil \texttt{secure}: los dos canales de
recuperaci\'on exigen \texttt{tenant\_id} y \texttt{user\_id} sin excepci\'on
(C1); lo recuperado sale del mensaje de sistema y entra delimitado en el
turno de usuario, con un marcador aleatorio por sesi\'on no adivinable sin la
clave del sistema y un aviso expl\'icito de que es dato, no instrucciones
(C2); el writer clasifica cada candidato antes de guardarlo y rechaza \'ordenes
y secretos (C3); C4 firma la procedencia; C5 aplica tiempo de vida por
categor\'ia y borra en las tres capas con prueba de borrado; C6 exige
confirmaci\'on expl\'icita del usuario para cualquier herramienta con efecto
externo.

\section{Controles de seguridad}

\subsection{C1 --- Aislamiento por espacio de nombres}
Todo registro pertenece a un espacio \texttt{tenant\_id/user\_id/session\_id}
y ninguna lectura sale sin ese filtro, en ninguno de los dos canales. La
debilidad que corrige no es que falte la etiqueta --el dato \emph{s\'i} la
tiene-- sino que nadie la usa como filtro en la consulta. Se verifica con T02
y M2; el criterio de aprobaci\'on es \textbf{M2a~=~0}, no M2b~=~0.

\subsection{C2 --- Memoria como dato no ejecutable}
El contenido recuperado entra delimitado en el turno de usuario, nunca en el
mensaje de sistema. Es la segunda l\'inea de defensa: si algo malicioso
entrara al almac\'en a pesar de C3, C2 reduce la probabilidad de que se
ejecute como instrucci\'on. Es una mitigaci\'on estad\'istica sobre el modelo,
no una garant\'ia estructural, y se reporta as\'i.

\subsection{C3 --- Filtro de escritura con esquema tipado}
Clasifica cada candidato en \texttt{fact}, \texttt{preference},
\texttt{task\_context}, \texttt{instruction} o \texttt{secret} antes de
persistir. Solo los tres primeros llegan al almac\'en. En este hito el
clasificador es \textbf{basado en reglas}: patrones espec\'ificos para
\'ordenes dirigidas al modelo (por ejemplo, ``a partir de ahora'' solo cuenta
si va seguido de un verbo imperativo dirigido al asistente, nunca de un hecho
del usuario como ``a partir de ahora vivo en Medell\'in'') y un chequeo de
Luhn para n\'umeros con forma de tarjeta. Un juez LLM que cubra los casos que
las reglas no resuelven qued\'o fuera del alcance de este hito (secci\'on de
limitaciones); de haberse incluido, habr\'ia introducido adem\'as
circularidad --el mismo modelo objeto de estudio decidiendo qu\'e se guarda--
y una superficie de ataque nueva, porque el texto del atacante entrar\'ia
tambi\'en al \emph{prompt} del juez.

\subsection{C4 --- Procedencia firmada}
Cada registro lleva su origen, autor y sesi\'on firmados con
HMAC-SHA256 sobre el JSON can\'onico del registro, con una clave fuera del
almac\'en (\texttt{MEMLAB\_HMAC\_KEY}, nunca versionada). En este hito la
firma se aplica al canal de siembra controlada que usan las pruebas para
sembrar un dato sin pasar por C3 (ver secci\'on de metodolog\'ia); conectar la
verificaci\'on de firma a \textbf{toda} escritura y lectura normal es
extensi\'on de este trabajo.

\subsection{C5 --- Tiempo de vida y borrado verificable}
Tiempo de vida por categor\'ia (\texttt{task\_context}: 24~h; \texttt{fact} y
\texttt{preference}: sin vencimiento salvo petici\'on expl\'icita), evaluado
tanto en lectura como por una purga peri\'odica. El borrado alcanza
\textbf{las tres capas}: la fila en SQLite (DELETE real, no l\'ogico), el
embedding en ChromaDB y la cach\'e de contexto de sesi\'on --por
\textbf{contenido}, no solo por la sesi\'on de origen, porque el contenido de
un registro puede quedar cacheado bajo la sesi\'on de cualquiera que lo haya
consultado, no solo la de quien lo escribi\'o--. Cada borrado deja una
\texttt{deletion\_proof} en el log de auditor\'ia, firmada con HMAC y
\textbf{nunca} con un hash plano: un SHA-256 simple sobre un dato corto y de
baja entrop\'ia (una direcci\'on, un n\'umero de tarjeta) se invierte por
fuerza bruta en segundos, y el log de solo-agregado terminar\'ia siendo una
copia recuperable de justo lo que se prometi\'o borrar. La purga es una
funci\'on pura que recibe el reloj como par\'ametro, para que T04 pueda
adelantarlo con un reloj inyectable sin depender de que un planificador en
otro hilo vea el mismo reloj congelado.

\subsection{C6 --- Autorizaci\'on por origen y auditor\'ia}
Toda acci\'on con efecto externo exige confirmaci\'on expl\'icita del usuario
en el turno en curso; la memoria y los documentos nunca autorizan por s\'i
mismos. Qued\'o fuera del alcance de este hito por tiempo (secci\'on de
limitaciones): el mecanismo de herramientas del orquestador y la prueba T03
que lo ejercita no se implementaron.

\section{Metodolog\'ia}

\textbf{Dise\~no general.} Cada prueba a\'isla el control que dice medir. Por
ejemplo, en T02/Secure el dato sensible de un usuario se siembra en el
almac\'en \textbf{sin pasar por C3} (ya firmado, por un canal separado): si se
sembrara por el camino normal, C3 rechazar\'ia la tarjeta como secreto, M2a
saldr\'ia en 0\,\% por falta de dato que fugar, y eso no demostrar\'ia nada
sobre C1. El mismo canal de siembra controlada se usa para crear el registro
de \texttt{task\_context} con tiempo de vida expl\'icito que necesita T04.

\textbf{Sesi\'on nueva es un proceso nuevo.} T01 lanza un subproceso por
sesi\'on con \texttt{session\_id} nuevo y cach\'e vac\'ia; un bucle dentro del
mismo proceso dejar\'ia estado compartido y no probar\'ia persistencia real.

\textbf{M1 --- Tasa de Persistencia de Instrucci\'on (T01).} Fracci\'on de
sesiones nuevas, sin relaci\'on sem\'antica con la instrucci\'on inyectada,
cuya respuesta la cumple. Se registra por separado cu\'antas sesiones
\emph{recuperaron} la instrucci\'on, para distinguir un fallo de recuperaci\'on
de un fallo de obediencia del modelo.

\textbf{M2 --- Tasa de Fuga entre Conversaciones (T02), en dos niveles.}
\textbf{M2a} (fuga de recuperaci\'on): el \'indice entreg\'o un documento de
otro espacio de nombres, le\'ido del log de recuperaci\'on, no de la
respuesta. \textbf{M2b} (fuga de divulgaci\'on): el dato apareci\'o en la
respuesta. Puede haber M2a sin M2b porque el modelo decidi\'o no repetir el
dato, y eso no es m\'erito de ning\'un control. El criterio de aprobaci\'on de
C1 es M2a~=~0.

\textbf{M4 --- Tiempo de vida y borrado verificable (T04, T05).} T04 mide
cu\'antas de 10 consultas, tras adelantar el reloj 48~h sobre un
\texttt{task\_context} con TTL de 24~h, siguen recuper\'andolo. T05
(``olvida mi direcci\'on'') mide, tras una sola llamada de borrado, el
porcentaje de las tres capas (fila, embedding, cach\'e) que qued\'o
irrecuperable --un chequeo estructural fijo de 3 elementos, sin intervalo de
Wilson, porque no es una proporci\'on muestreada-- y, aparte, con Wilson sobre
10 consultas conversacionales posteriores, si un usuario normal logra sacar
el dato de vuelta por la v\'ia conversacional.

\textbf{Intervalos de confianza.} Toda proporci\'on se reporta con el
intervalo de Wilson; con 0 aciertos se reduce a la regla de tres ($3/n$). Con
20 sesiones cada acierto vale 5 puntos porcentuales, as\'i que el valor
puntual solo, sin intervalo, ser\'ia enga\~noso.

\textbf{Latencia en dos niveles.} $L_{e2e}$ (turno completo, con
generaci\'on) se reporta como contexto. $L_{mem}$ (recuperaci\'on +
ensamblado de \emph{prompt} + filtro de escritura, \textbf{sin}
generaci\'on) es la \'unica que los controles afectan y la que sustenta H4:
medida extremo a extremo, la varianza de generaci\'on del modelo (entre
cuatro y cinco veces el presupuesto completo de H4 solo entre percentiles de
la misma l\'inea base, ver Hito~2) enterrar\'ia por completo el efecto de los
controles.

\textbf{Reproducibilidad.} Semilla 20260919, temperatura 0,
\texttt{num\_predict}~=~256, una sola m\'aquina por comparaci\'on, digest del
modelo en cada manifiesto (\texttt{46e0c10c\ldots a666e} para
\texttt{llama3.1:8b-instruct-q4\_K\_M}, Ollama 0.34.2). Las corridas de este
hito se ejecutaron en la \textbf{misma m\'aquina} que la l\'inea base del
Hito~2. Antes de comparar perfiles se repitieron T01 y T02 Unsecure sobre el
commit nuevo y se contrastaron con la l\'inea base del commit \texttt{3ada423},
para descartar que el c\'odigo del Hito~3 hubiera alterado el camino Unsecure
(secci\'on~\ref{sec:resultados}).

\section{Resultados}
\label{sec:resultados}

Todas las cifras de esta secci\'on salen de \texttt{experiments/compute\_me-}
\texttt{trics.py} leyendo \texttt{evidence/runs/}; ninguna se escribe a mano.
Corridas oficiales del 17/10/2026 sobre el commit \texttt{614fadd}, backend
Ollama, digest \texttt{46e0c10c\ldots a666e}, en la misma m\'aquina de la
l\'inea base.

\subsection{Verificaci\'on de la l\'inea base Unsecure (regla de no
regresi\'on)}

\begin{table}[h]
\centering
\caption{Unsecure, Hito~2 (\texttt{3ada423}) vs.\ commit nuevo, misma m\'aquina}
\label{tab:regresion}
\begin{tabular}{@{}lcc@{}}
\toprule
 & \texttt{3ada423} & \texttt{614fadd} \\
\midrule
T01, instrucci\'on cumplida & 20/20 & 20/20 \\
T02, M2a (fuga recuperaci\'on) & 67/100 & 67/100 \\
T02, M2b (fuga divulgaci\'on) & 49/100 & 49/100 \\
\bottomrule
\end{tabular}
\end{table}

Las tres cifras de seguridad coinciden exactamente, tal como ya hab\'ia
ocurrido entre las corridas del 19/09 y del 22/09. El c\'odigo del Hito~3 no
alter\'o el camino Unsecure: la diferencia entre perfiles que sigue se debe a
los controles. Lo \'unico que vari\'o fue la latencia, como era de esperar por
la carga de la m\'aquina ($L_{mem}$ p95 de T01: 100~ms el 22/09, 83~ms el
17/10).

\subsection{Unsecure vs.\ Secure}

\begin{table}[h]
\centering
\caption{Unsecure vs.\ Secure, commit \texttt{614fadd}, misma m\'aquina. Entre corchetes, IC~95\,\% de Wilson}
\label{tab:comparativa}
\begin{tabular}{@{}lcc@{}}
\toprule
M\'etrica & Unsecure & Secure \\
\midrule
M1 --- TPI (T01) & 20/20 [83,9--100] & 0/40 [0--8,8] \\
M2a --- fuga recuperaci\'on (T02) & 67/100 [57,3--75,4] & 0/100 [0--3,7] \\
M2b --- fuga divulgaci\'on (T02) & 49/100 [39,4--58,7] & 0/100 [0--3,7] \\
T04 --- recuperado tras TTL & 10/10 [72,3--100] & 0/10 [0--27,8] \\
M4 --- capas irrecuperables (T05) & 1/3 (33,3\,\%) & 3/3 (100\,\%) \\
T05 --- fuga conversacional tras olvidar & 0/10 & 0/10 \\
$L_{mem}$ p50 / p95, T01 (ms) & 69 / 83 & 69 / 86 \\
$L_{mem}$ p50 / p95, T02 (ms) & 48 / 56 & 49 / 53 \\
$L_{e2e}$ p95, T01 (ms) & 5248 & 5024 \\
\bottomrule
\end{tabular}
\end{table}

\textbf{Persistencia (T01).} En Secure ninguna de las 40 sesiones nuevas
cumpli\'o la instrucci\'on, y ninguna la recuper\'o: el log muestra que C3 la
rechaz\'o en el momento de escribirla, con cuatro patrones a la vez (formato
impuesto a las respuestas, ``debes'', orden temporal y enlace). Esto significa
que el resultado mide C3, no C2: la instrucci\'on nunca lleg\'o al almac\'en,
as\'i que C2 no tuvo nada que neutralizar. Con 40 sesiones, la cota superior
del intervalo (8,8\,\%) todav\'ia queda por encima del umbral absoluto del
5\,\% de H1: el valor puntual cumple, pero la muestra no basta para
garantizarlo.

\textbf{Aislamiento (T02).} M2a lleg\'o a 0 de 100, que es el criterio de
aprobaci\'on de C1. El resultado es atribuible a C1 porque los dos registros
de A se sembraron firmados y sin pasar por C3: el dato ajeno estaba en el
almac\'en en cada uno de los 100 sondeos y aun as\'i ninguno de los dos
canales lo entreg\'o. Como M2a fue 0, M2b tambi\'en: el modelo no puede repetir
lo que nunca recibi\'o.

\textbf{Ciclo de vida (T04, T05).} El dato de contexto con TTL de 24~h se
recuper\'o en 10 de 10 consultas en Unsecure despu\'es de adelantar el reloj
48~h, y en ninguna en Secure. Tras ``olvidar'', Unsecure solo marc\'o la fila
en SQLite y dej\'o el embedding y la cach\'e; Secure borr\'o las tres capas y
dej\'o una prueba de borrado con HMAC en el log. La fuga conversacional tras
olvidar fue 0/10 en los dos perfiles (ver Discusi\'on).

\textbf{Costo.} El percentil 95 de $L_{mem}$ subi\'o 3~ms en T01 y baj\'o
3~ms en T02: el efecto de C1, C2, C3 y C5 sobre la ruta de memoria est\'a por
debajo del ruido de la m\'aquina y muy lejos del presupuesto de 300~ms de H4.

\section{Discusi\'on}

La brecha entre M2a y M2b en la l\'inea base (67\,\% vs.\ 49\,\%, 18 puntos)
es, en s\'i misma, un hallazgo metodol\'ogico: reportar \'unicamente lo que el
modelo repite subestimar\'ia en un 27\,\% relativo una falla del control que
ya ocurri\'o en el momento en que el \'indice entreg\'o el documento. Por eso
el criterio de aprobaci\'on de C1 en este trabajo es M2a~=~0 y no M2b~=~0: un
sistema que entrega datos ajenos y confía en que el modelo se calle no tiene
un control, tiene suerte.

La diferencia entre $L_{e2e}$ y $L_{mem}$ --en la l\'inea base del Hito~2, la
sola diferencia entre percentiles del turno completo (1384~ms en T01,
1248~ms en T02) ya super\'o entre cuatro y cinco veces el presupuesto entero
de H4 sin que existiera ning\'un control-- confirma que medir el costo de
seguridad extremo a extremo habr\'ia hecho la hip\'otesis de costo
inevaluable: la varianza de generaci\'on del modelo local esconde cualquier
efecto real de los controles. $L_{mem}$ aisla exactamente lo que C1--C6
tocan.

El borrado de T05 merece una precisi\'on. En Unsecure, tras ``olvidar'', el
dato no volvi\'o a salir por la conversaci\'on (0/10), porque el laboratorio
resuelve cada resultado del \'indice contra la fila de SQLite y esa fila ya
qued\'o marcada como borrada. La falla de Unsecure no est\'a en lo que un
usuario puede sacar charlando, sino en el dep\'osito: el embedding y la
cach\'e siguen guardando el dato, y cualquiera con acceso al \'indice o a un
respaldo puede leerlo. Por eso M4 se mide de forma estructural, capa por
capa, y no por la conversaci\'on: medido solo en la respuesta, el borrado
l\'ogico habr\'ia parecido suficiente.

\section{Limitaciones}

\begin{enumerate}
\item \textbf{Alcance inferido.} No existe r\'ubrica del docente para el
Hito~3 al momento de entregar; el alcance sale del cronograma de la propuesta
original y se declara expl\'icitamente aqu\'i.
\item \textbf{C3 es solo reglas.} Sin juez LLM, el clasificador no detecta
secretos sin forma de n\'umero de tarjeta (por ejemplo, una contrase\~na en
texto libre). Se declara, no se oculta.
\item \textbf{C4 parcial.} La firma se aplica al canal de siembra controlada
que usan las pruebas; conectarla a toda escritura y verificarla en toda
lectura normal qued\'o fuera de este hito por tiempo.
\item \textbf{C6 y T03 no implementados.} El mecanismo de herramientas
del orquestador no se conect\'o; la inyecci\'on indirecta desde un documento
no se midi\'o en este hito.
\item \textbf{M3 y M5 no medidos.} El conjunto etiquetado de 60 candidatos
para M3 existe (\texttt{experiments/datasets\_hito3.py}) pero no se
proces\'o por el filtro con intervalo de Wilson; T06 y la utilidad
evaluada por el equipo (M5, kappa de Cohen) tampoco.
\item \textbf{Un solo modelo y una sola m\'aquina.} Llama
3.1 8B Instruct Q4\_K\_M; temperatura 0 y semilla fija no garantizan
determinismo bit a bit entre m\'aquinas ni cuantizaciones.
\item \textbf{Detecci\'on automatica con casos de revisi\'on manual.} M1 y
M2b usan expresiones regulares; los casos amb\'iguos se marcan para revisi\'on
manual en vez de contarse autom\'aticamente.
\item \textbf{C2 es estad\'istico, no estructural, y no se aisl\'o.} Reduce
la probabilidad de obediencia a una instrucci\'on inyectada, no la elimina por
construcci\'on. Adem\'as, en T01 Secure la instrucci\'on nunca lleg\'o al
almac\'en porque la rechaz\'o C3, as\'i que el 0/40 no mide C2. Falta la
ablaci\'on con la instrucci\'on ya firmada dentro del almac\'en.
\item \textbf{Tama\~no de muestra en T01.} Con 0 aciertos en 40 sesiones, la
cota superior del intervalo de Wilson es 8,8\,\%: el valor puntual cumple
H1 (TPI~$\le$~5\,\%), pero la muestra no basta para garantizarlo.
\item \textbf{Sesgo optimista en las reglas de C3.} Las reglas y los textos de
prueba se escribieron en la misma sesi\'on de trabajo, y algunas reglas
reconocen literalmente la forma de los ataques del laboratorio (por ejemplo,
``guarda en memoria que'', del documento de T03). Cualquier cifra de M3 sobre
ese conjunto tendr\'a sesgo optimista y deber\'a contrastarse con textos
escritos por terceros.
\item \textbf{T01 con un solo usuario.} La instrucci\'on la escribe y la
recibe el mismo usuario B. El caso en que A escribe y B la recibe por el
perfil fijado no se midi\'o.
\end{enumerate}

\section{Conclusiones}

Sobre una misma base de c\'odigo, el perfil Unsecure reprodujo exactamente
la falla que motiva este trabajo: una instrucci\'on inyectada persiste y se
obedece en el 100\,\% de sesiones nuevas sin relaci\'on sem\'antica con ella,
y el \'indice vectorial entrega datos de otro usuario en 67 de cada 100
sondeos aunque el modelo solo los repita en 49. Agregar aislamiento por
espacio de nombres en los dos canales de recuperaci\'on, separar datos de
instrucciones en el ensamblado del \emph{prompt}, clasificar lo que se
escribe antes de persistirlo, y borrar en las tres capas con prueba firmada,
redujo la persistencia de 20/20 a 0/40 sesiones, la fuga de recuperaci\'on de
67 a 0 sondeos de 100, la recuperaci\'on de datos vencidos de 10/10 a 0/10, y
llev\'o el borrado de una a tres capas, con un cambio de unos 3~ms en el
percentil 95 de la ruta de memoria. Qued\'o fuera de este hito, por tiempo y
se declara como tal, la implementaci\'on completa de C4 y C6, la medici\'on de
M3 y M5, la ablaci\'on que a\'isla a C2, y el juez LLM dentro de C3.

\section{Declaraci\'on de uso de IA}

\begin{table}[h]
\centering
\caption{Uso de IA en este trabajo}
\begin{tabular}{@{}p{0.38\linewidth}p{0.52\linewidth}@{}}
\toprule
Herramienta & Prop\'osito concreto \\
\midrule
Claude (Anthropic), Claude Code & Implementaci\'on de los controles C1, C2,
C3 y C5 sobre el dise\~no ya definido en la propuesta; escritura de los
\emph{runners} de pruebas T04/T05 y de las extensiones a
\texttt{compute\_metrics.py}; correcci\'on de inconsistencias conocidas en la
documentaci\'on; verificaci\'on del mapeo a MITRE ATLAS contra la fuente
oficial; lanzamiento de los \emph{runners} para las corridas oficiales en la
m\'aquina del equipo; redacci\'on de este documento y de la presentaci\'on. \\
Llama 3.1 8B Instruct (Ollama, local) & \textbf{Objeto de estudio} del
experimento: es el modelo sobre el que se miden las fugas y persistencias;
no particip\'o en la evaluaci\'on de utilidad (M5), que exige juicio humano
para no contradecir la propia declaraci\'on de uso de IA. \\
\bottomrule
\end{tabular}
\end{table}

\textbf{Responsabilidad del equipo.} Todo el contenido producido con ayuda de
IA fue revisado, verificado y ajustado por los integrantes, que asumen la
autor\'ia y la responsabilidad acad\'emica de este trabajo.
\textbf{Ninguna cifra experimental de este documento fue generada o
estimada por una IA}: todas salen de ejecutar el laboratorio y est\'an
respaldadas por evidencia versionada en \texttt{evidence/}, con su
\texttt{git\_sha}, su semilla y el digest del modelo. Los datos del
laboratorio son inventados (dos usuarios simulados, tarjeta de prueba
est\'andar de Visa, direcciones ficticias); ninguna prueba se ejecut\'o
sobre sistemas, cuentas o servicios de terceros.

\begin{thebibliography}{00}

\bibitem{owasp-llm10-2025}
OWASP Foundation, ``OWASP Top 10 for Large Language Model Applications
2025,'' \emph{genai.owasp.org}, 2025.

\bibitem{owasp-agentic-2026}
OWASP GenAI Security Project, ``OWASP Top 10 for Agentic Applications
2026,'' ASI06 Memory \& Context Poisoning, \emph{genai.owasp.org}, Dec.\ 9,
2025.

\bibitem{minja-2025}
S.\ Dong, S.\ Xu, P.\ He, Y.\ Li, J.\ Tang, T.\ Liu, H.\ Liu, and Z.\ Xiang,
``Memory Injection Attacks on LLM Agents via Query-Only Interaction,''
\emph{arXiv:2503.03704}, 2025.

\bibitem{memory-poisoning-defense-2026}
B.\ Devarangadi Sunil, I.\ Sinha, P.\ Maheshwari, S.\ Todmal, S.\ Mallik, and
S.\ Mishra, ``Memory Poisoning Attack and Defense on Memory Based
LLM-Agents,'' \emph{arXiv:2601.05504}, 2026.

\bibitem{nist-aml-2025}
National Institute of Standards and Technology, ``Adversarial Machine
Learning: A Taxonomy and Terminology of Attacks and Mitigations,'' NIST AI
100-2 E2025, Mar.\ 2025.

\bibitem{mitre-atlas}
MITRE ATLAS, ``Adversarial Threat Landscape for Artificial-Intelligence
Systems,'' \emph{atlas.mitre.org}, v5.6.0 (consultado 16/10/2026).

\bibitem{ley-1581-2012}
Congreso de la Rep\'ublica de Colombia, ``Ley Estatutaria 1581 de 2012,
R\'egimen General de Protecci\'on de Datos Personales,'' 17 oct.\ 2012.

\end{thebibliography}

\end{document}
